\chapter{Frameworks for Optimal Control Problem Formulation}
Process of solving optimal control problem has 3 parts. Problem defintion,
solving of nonlinear program, optional mesh refinment

\begin{figure}[h]
  \centering
  \begin{tikzpicture}[
      node distance=1.4cm,
      block/.style={rectangle, draw, thick, font=\small, minimum width=4cm, minimum height=0.9cm, align=center, fill=white},
      exit/.style={ellipse, draw, thick, font=\small, minimum width=3cm, minimum height=0.8cm, align=center, fill=white}
    ]
    \node[block] (def) {1. Problem definition};
    \node[block, below=of def] (nlp) {2. NLP solve};
    \node[block, below=of nlp] (mesh) {3. Mesh refinement (optional)};
    \node[exit, below=of mesh] (sol) {Solution};
    \draw[->, thick] (def) -- (nlp);
    \draw[->, thick] (nlp) -- (mesh);
    \draw[->, thick] (mesh) -- node[right, font=\small]{converged} (sol);
    \draw[->, thick] let \p1=(mesh.east), \p2=(nlp.east) in
    (mesh.east) -- (\x1+1.5cm,\y1) -- node[right, font=\small, align=left]{needs\\ refinement} (\x1+1.5cm,\y2) -- (nlp.east);
  \end{tikzpicture}
  \caption{Optimal control problem solving pipeline.}
  \label{fig:ocp-pipeline}
\end{figure}

The concensus in optimal control community about tools for prolbem formulation
is that:
\begin{quote}
  There is no shortage of optimal control software based on pseudospectral approaches. There are a number of other optimal control libraries that tackle similar kinds of problems, such as OTIS4, GPOPS-II,and CASADI.
\end{quote}
\cite{Falck2021}\cite{DymosDocs}
This chapter is aobut tool that define the OCP, and supply derivatives of the problem, they cannot solve the problem by themselves.
\section{GPOPS-II}
GPOPS-II \cite{pattersonGPOPSIIMATLABSoftware2014} is perhaps the most advaced
tool for otimal control, it was developed ad University of Florida by Michael
patterson and Anil V. Rao. other tools allows defininng the problem, but
GPOPS-II is a level higher. It also includes advanced mesh refinement
strategies, multiple phase OCP and automatic scaling. For solving the problem
it uses IPOPT or SNOPT. It interaces with MATLAB and C++, uses ADigator
\cite{weinstein2017algorithm} for automatic differetnation. This tool is
commercial and closed source. Implementaion in this thesis is heavily based on
work of community from University of Florida.
\section{CasADi}
CasADi is an open-source tool for nonlinear optimization and algorithmic
differentiation.\cite{CasADi}\cite{anderssonCasADiSoftwareFramework2019}.
interfaces with C++,Python,MATLAB. It provides a framework INE SLOVO to build
the problem, and calculate the hessian and jacobian for nlp solver it does not
have any mesh refinemnt algorithms. Commonly paired with IPOPT for nonliner
programming, but is compatible with large variety of solvers.

\section{dymos + OpenMDAO}
Python has hp mesh refinemnt uses OpenMDAO \cite{Gray2019a} developed at NASA
Glenn research center for evaluating aircraft concepts. Lobatto and radau
collocation as default \cite{Falck2021}. According to forum post from julia
discourse, which may be biased towards julia it is slower compared to infiteOpt
in julia \cite{DoesAnybodyKnow2024}. But it has mesh refinemnt so that is good.
\section{Acados}

\cite{Verschueren2021}

\section{Julia + JuMP/ infititeopt/ examomdels}
ExaModels, JuMP

JuMP is the go to tool for algebraic modelling in julia, it is mature,
versatile. Supports Linear programming, quadartic programming, mixed linear
programming, Nonlinear programming. ExaModels \cite{shin2023accelerating} is a
specialized algebraic modelling tool, it utilisies Single instruction, multiple
data methods,enabling derivative evaluation on GPU accelerators. Coupled with
MadNLP solver, the entire computation of nonlinear optimisation can be run on a
GPU.

Julia options
\begin{enumerate}
  \item Julia + jump
  \item Julia + jump + infinte opt
  \item julia + ExaModels
\end{enumerate}

Infite opt already has various collocation methods implemented. For learning
purposes and in fear that in future down the road i would realise that
infniteopt blocks all the stuff i want to do i decided to use jump, without
infite opt.

No julia package has  mesh refinement
I chose julia + jump,


\begin{table}[h]
  \centering
  \small
  \begin{tabular}{|p{4.5cm}|p{4.5cm}|p{4.5cm}|}
    \hline
    \textbf{Aspect} & \textbf{Julia + JuMP + InfiniteOpt} & \textbf{Python + OpenMDAO + dymos} \\
    \hline
    Language / runtime
      & Julia, JIT-compiled, no GIL
      & Python, interpreted, GIL (NumPy releases) \\
    \hline
    License
      & MIT (all components)
      & Apache 2.0 \\
    \hline
    Algebraic modelling style
      & Macro-based (\texttt{@variable}, \texttt{@constraint}, \texttt{@objective})
      & Component graph (\texttt{ExplicitComponent}, \texttt{Group}, promotes/connects) \\
    \hline
    Derivatives
      & AD via MathOptInterface; sparsity from expression graph
      & MAUD unified derivatives; mix analytic / AD / FD / complex-step per component, framework chains them \\
    \hline
    Transcription methods
      & Finite difference (Forward, Central, Backward); orthogonal collocation on finite elements with Gauss--Lobatto nodes only
      & Radau pseudospectral, Gauss--Lobatto, explicit shooting, ExplicitShooting RK \\
    \hline
    Mesh / grid refinement
      & None built-in; user must re-solve on refined supports
      & Built-in ph- and hp-adaptive grid refinement \\
    \hline
    Multi-phase trajectories
      & Manual; user concatenate \texttt{InfiniteModel}s and link
      & Native \texttt{Trajectory} and \texttt{Phase} with linkage constraints \\
    \hline
    Multidisciplinary coupling
      & Not supported; pure infinite-dim model
      & Native MDO coupling with implicit/explicit components, cyclic groups, nonlinear/linear sub-solvers \\
    \hline
    NLP solvers
      & IPOPT, MadNLP, Knitro, SNOPT and any MOI-compatible solver
      & IPOPT, SNOPT, scipy via pyOptSparse \\
    \hline
    GPU acceleration
      & Yes via ExaModels + MadNLP (SIMD on GPU)
      & Not native \\
    \hline
    Parallelism
      & Native multithreading, Distributed, MPI
      & MPI across components \\
    \hline
    Sparse total Jacobian
      & Yes, from MOI expression graph
      & Yes, MAUD assembles sparse total Jacobian from component partials \\
    \hline
    Maturity / community
      & JuMP mature in OR / energy; InfiniteOpt newer, smaller user base
      & NASA-backed, large aerospace user base \\
    \hline
    Learning curve for OCP only
      & Low; close to mathematical formulation
      & Higher; OpenMDAO component model required even for single-discipline OCP \\
    \hline
  \end{tabular}
  \caption{Comparison of Julia (JuMP + InfiniteOpt) and Python (OpenMDAO + dymos) stacks for optimal control problem formulation.}
  \label{tab:julia-vs-python-ocp}
\end{table}

